% SEGMENT OLP-0699-S01
% Part: proof-theory
% Chapter: sequent-calculus
% Section: interpretation-rules

\documentclass[../../../include/open-logic-section]{subfiles}

\begin{document}

\olfileid{pt}{seq}{irl}

\olsection{விதிகளின் பொருள்விளக்கம்}

ஒரு தொடரணியின் பொருள்விளக்கத்தை நினைவுகூர்க:
தொடரணி $\Gamma \Sequent \Delta$, !!a{structure}~$\Struct{M}$ இல்
மெய்யாக இருப்பது, அதன் இடப்புறத்தில் உள்ள
குறைந்தது ஒரு $!A \in \Gamma$ பொய்யாகவோ,
அதன் வலப்புறத்தில் உள்ள குறைந்தது ஒரு
$!A \in \Delta$ மெய்யாகவோ இருந்தால்
மட்டுமே.
\Log{G3c} இன் தருக்க விதிகளை, அவற்றின் முதன்மை
!!{formula}s இன் மெய்மை நிபந்தனைகளை
விதிவடிவமாக்குவனவாகப் படிக்கலாம்.
\RightR{\lif} ஐ எடுத்துக்கொள்க. அதன் முதன்மை
!!{formula} $!A \lif !B$ ஆகும். $!A$ பொய்யாக
இருந்தாலோ $!B$ மெய்யாக இருந்தாலோ அது
மெய்யாகும். எனவே $\ \Sequent !A \lif !B$
மெய்யாக இருப்பதும் $!A \Sequent !B$
மெய்யாக இருப்பதும் ஒரே நிபந்தனை.
இவ்விரு தொடரணிகளின் இடப்புறங்களிலும்
வலப்புறங்களிலும் முறையே சூழல் !!{formula}s
ஆன $\Gamma$, $\Delta$ வைச் சேர்த்தால்,
\RightR{\lif} விதியின் முடிவும் முன்கூற்றும்
துல்லியமாகக் கிடைக்கும்.
\LeftR{\lif} க்கும் இதேபோன்ற நிலை:
$!A$ மெய்யாகவும் $!B$ பொய்யாகவும் இருந்தால்
$!A \lif !B$ பொய்யாகும். ஆகவே
$!A \lif !B \Sequent \ $ மெய்யாக இருப்பதும்
$\ \Sequent !A$ மற்றும் $!B \Sequent \ $
இரண்டும் மெய்யாக இருப்பதும் ஒரே நிபந்தனை.
\LeftR{\lif} இன் முடிவு, அதன்
முன்கூற்றுகளின் இணையலுக்குச் சமானமானது.
இந்த அமைப்பு \Log{G3c} விதிகளின்
சரித்தன்மையை உறுதிப்படுத்துகிறது:
எல்லா முன்கூற்றுகளும் மெய்யென்றால்
முடிவும் மெய். அவை நிறைவுத்தன்மையையும்
உறுதிப்படுத்துகின்றன.

% SEGMENT OLP-0699-S02
உண்மையில், மெய்ம்மதிப்புச் சார்ந்த எந்த
இணைப்பியின் மெய்மை நிபந்தனைகளையும்
இவ்வாறு தொடரணிகளின் தொகுப்பாக விவரிக்கலாம்.
பாரம்பரியத் தருக்கத்தின் ஒவ்வொரு மெய்ம்மதிப்புச்
சார்பையும் இணைவு இயல்புவடிவிலுள்ள
!!a{formula} ஆல் விவரிக்க முடிவதே இதற்குக்
காரணம்: அது அணு !!{formula}s மற்றும்
அவற்றின் மறுப்புகளின் பிரிப்பிணைவுகளை
இணைத்த வாய்பாடு. எடுத்துக்காட்டாக,
$!A \oplus !B$ என்பது $!A$, $!B$
ஆகிய இரண்டில் சரியாக ஒன்று மட்டும்
மெய்யாக இருக்கும்போது மெய்யாகும் என்று
கொள்க; இதுவே “பிரத்தியேக அல்லது”.
அப்போது $!A \oplus !B$ மெய்யாக
இருப்பதும் $(!A \lor !B) \land
(\lnot !A \lor \lnot !B)$ மெய்யாக
இருப்பதும் ஒரே நிபந்தனை. மேலும்
$(\lnot !A \lor !B) \land
(!A \lor \lnot !B)$ மெய்யாக இருந்தால்
அந்த இணைப்பி பொய்யாகும்.
எனவே $\oplus$ க்கான தொடரணிக்
கணிய விதிகளை இவ்வாறு அறிமுகப்படுத்தலாம்:
\begin{align*}
&
\Axiom$\Gamma \fCenter \Delta, !A, !B$
\Axiom$!A, !B, \Gamma \fCenter \Delta$
\RightLabel{\RightR{\oplus}}
\BinaryInf$\Gamma \fCenter \Delta, !A \oplus !B$
\DisplayProof
\\
& \Axiom$!A, \Gamma \fCenter \Delta, !B$
\Axiom$!B, \Gamma \fCenter \Delta, !A$
\RightLabel{\LeftR{\oplus}}
\BinaryInf$!A \oplus !B, \Gamma \fCenter \Delta$
\DisplayProof
\end{align*}
இவ்விதிகளும் சரியானவையும்
நிறைவுத்தன்மையுடையவையும் ஆகும்.
அவற்றின் முடிவுகள் மெய்யாக இருப்பதும்
அவற்றின் எல்லா முன்கூற்றுகளும்
மெய்யாக இருப்பதும் ஒரே நிபந்தனை.

% SEGMENT OLP-0699-S03
\begin{prob}
$!A \downarrow !B$ என்பது $!A$ உம்
$!B$ உம் இரண்டுமே மெய்யல்லாதபோது
மட்டுமே மெய்யாகும் (“இரண்டுமன்று”, NOR) என்று
கொள்க. அதற்கான \Log{G3c} விதிகள் யாவை?
($!A \downarrow !B$ மெய்யாக இருக்கும்போது
மட்டுமே ஒரே நேரத்தில் மெய்யாகும்
இரு தொடரணிகளைக் கண்டறிக. அவை
\RightR{\downarrow} விதியைத் தரும்.
பிறகு $!A \downarrow !B$ பொய்யாக
இருக்கும்போது மட்டுமே மெய்யாகும்
ஒரு தொடரணியைக் கண்டறிக. அது
\LeftR{\downarrow} விதியைத் தரும்.)
\end{prob}

% SEGMENT OLP-0699-S04
\Log{G3c} இல் ஒவ்வோர் இணைப்பிக்கும்
அளவையடைக்கும் சரியாக இரண்டு விதிகள்
உள்ளன: ஓர் இட விதியும் ஒரு வல விதியும்.
ஆனால் \Log{G1c} இல் இரண்டு
\LeftR{\land} விதிகளும் இரண்டு
\RightR{\lor} விதிகளும் உள்ளன.
இதற்குக் காரணம், \Log{G1c} க்கு
மாதிரியாக விளங்கிய ஜென்ட்சனின்
மூல முறைமை~\Log{LK} க்கு,
உள்ளுணர்வுவாதத் தருக்கம் என்ற முக்கியமான
பாரம்பரியமல்லாத தருக்கத்திற்குரிய
\Log{LJ} என்ற மாறுவடிவமும் இருந்ததே.
உள்ளுணர்வுவாதத் தருக்கத்திற்குச் சரியானதும்
நிறைவானதுமான முறைமையைப் பெற,
\Log{LJ} இல் ஒவ்வொரு தொடரணியின்
வலப்புறமும் அதிகபட்சம் ஒரு
!!{formula} மட்டுமே கொண்டிருக்க
வேண்டும். இதனால் \Log{G3c} இன்
\RightR{\lor} விதியைப் பயன்படுத்த
முடியாது: அதன் முன்கூற்றின்
வலப்புறத்தில் $!A$, $!B$ இரண்டும்
உள்ளன. ஆகவே ஜென்ட்சன் \Log{LJ}
க்காக ஒரு சோடி \RightR{\lor}
விதிகளைப் பயன்படுத்தினார்;
\Log{LK} க்கும் அதே விதிகளைப்
பயன்படுத்தினார். இதேபோல்,
உள்ளுணர்வுவாத மாறுவடிவமான~\Log{G1i}
என்பது, எல்லாத் தொடரணிகளின்
வலப்புறங்களிலும் அதிகபட்சம் ஒரு
!!{formula} மட்டுமே இருக்க வேண்டும்
என்ற கட்டுப்பாட்டுடன் கூடிய~\Log{G1c}
ஆகும். (\Log{G3c} க்கும் ஓர்
உள்ளுணர்வுவாத மாறுவடிவம் உண்டு;
ஆனால் அதன் சில விதிகள்~\Log{G3c}
இன் ஒத்த விதிகளிலிருந்து வேறுபடுகின்றன.
எடுத்துக்காட்டாக, அதுவும் ஒரு சோடி
\RightR{\lor} விதிகளைப் பயன்படுத்துகிறது.)

% SEGMENT OLP-0699-S05
\Log{G1c} இன் கட்டமைப்பு விதிகள்
சரியானவை என்பதை எளிதில் காணலாம்.
எடுத்துக்காட்டாக, $\Gamma \Sequent
\Delta$ வின் இடப்புறத்தில் பொய்யான
!!{formula} ஒன்றோ வலப்புறத்தில்
மெய்யான !!{formula} ஒன்றோ
இருந்தால், $\Gamma \Sequent
\Delta, !A$ விலும் அதே
வாய்பாடு இருக்கும். $!A$
மெய்யா பொய்யா என்பது
பொருட்டல்ல. எனவே
\RightR{\Weakening} இன்
முன்கூற்று மெய்யென்றால் அதன்
முடிவும் மெய். இதன் மறுதிசை
பொதுவாக மெய்யல்ல என்பதை
கவனிக்க: $!A$ மெய்யாக
இருப்பதாலேயே முடிவு
மெய்யாகலாம்; அப்போது
முன்கூற்று மெய்யாகும்
என்ற உறுதி இல்லை.

% SEGMENT OLP-0699-S06
கட்டமைப்பு விதிகள் ஏன் தேவை?
எடுத்துக்காட்டாக, $\ \Sequent !A
\lor \lnot !A$ ஐ !!{prove}
செய்யச் சுருக்கல்
(\RightR{\Contraction},
\LeftR{\Contraction}) தேவைப்படுகிறது.
$!A \Sequent !A$ இலிருந்து
தொடங்கினால், முதலில்
\RightR{\lnot} மூலம்
$\Sequent !A, \lnot !A$
கிடைக்கும். பின்னர்
\RightR{\lor} இன் இரு
வடிவங்களையும் தொடர்ச்சியாகப்
பயன்படுத்தலாம்: முதலில்
$\lnot !A$ இலிருந்து
$!A \lor \lnot !A$ ஐயும்,
அடுத்து $!A$ இலிருந்து
$!A \lor \lnot !A$ ஐயும்
பெறலாம். கிடைப்பது
$\ \Sequent !A \lor \lnot !A,
!A \lor \lnot !A$ ஆகும்.
எனவே \Log{G1c} இல்
$\Sequent !A \lor \lnot !A$
ஐப் பெற, ஒரே !!{formula}
இருமுறை வருவதை !!a{proof}
இல் “சுருக்கும்” வசதி
தேவை:
\[
\Axiom$!A \fCenter !A$
\RightLabel{\RightR{\lnot}}
\UnaryInf$\fCenter !A, \lnot !A$
\RightLabel{\RightR{\lor}}
\UnaryInf$\fCenter !A, !A \lor \lnot !A$
\RightLabel{\RightR{\lor}}
\UnaryInf$\fCenter !A \lor \lnot !A, !A \lor \lnot !A$
\RightLabel{\RightR{\Contraction}}
\UnaryInf$\fCenter !A \lor \lnot !A$
\DisplayProof
\]

% SEGMENT OLP-0699-S07
உண்மையில், தளர்த்தலோ சுருக்கலோ
இல்லாமல் \Log{G1c} இல் !!{prove}
செய்ய முடியாத சரியான
தொடரணிகளும் உள்ளன.
எடுத்துக்காட்டாக, \Log{G1c}
இல் $!A \Sequent !B \lif !A$
என்பதன் !!a{proof},
\RightR{\lif} விதியில்தான்
முடிய வேண்டும்; ஏனெனில்
$!B \lif !A$ மட்டுமே
முதன்மை வாய்பாடாக இருக்க
முடியும். இதற்கு $!B, !A
\Sequent !A$ முன்கூற்றாகத்
தேவை. அது $!A \Sequent !A$
என்ற அடிக்கோளிலிருந்து
\LeftR{\Weakening} மூலம்
பெறப்பட வேண்டும்:
\[
\Axiom$!A \fCenter !A$
\RightLabel{\LeftR{\Weakening}}
\UnaryInf$!B, !A \fCenter !A$
\RightLabel{\RightR{\lif}}
\UnaryInf$!A \fCenter !B \lif !A$
\DisplayProof
\]

% SEGMENT OLP-0699-S08
\Log{G3c} இல் சுருக்கலும்
தளர்த்தலும் தனி விதிகளாகத்
தேவையில்லை. தளர்த்தல்
அடிக்கோள்களிலேயே பொதிந்துள்ளது.
இரு முன்கூற்றுகள் உள்ள விதிகளில்,
சூழல் $\Gamma$, $\Delta$ வின்
இரு நகல்களையும் ஒன்றாகப்
பகிர்வதன் மூலம் பெரும்பாலும்
சுருக்கலைத் தவிர்க்கலாம்.
\Log{G1c}, \Log{G3c}
இரண்டிலும் இத்தகைய
“சூழல் பகிரும்” விதிகள் உள்ளன.
ஒரு முன்கூற்று மட்டுமே உள்ள
விதிகளில், \RightR{\lor},
\LeftR{\land},
\RightR{\lexists},
\LeftR{\lforall} ஆகியவற்றுக்கே
சுருக்கல் தேவைப்படுகிறது.
\Log{G3c} இல் இவ்விதிகள்
கொள்ளும் வடிவங்கள் சுருக்கலை
முற்றிலும் தேவையற்றதாக்குகின்றன.
இரு முன்கூற்று விதிகளுக்கு
ஒன்றுக்கொன்று சாராத சூழல்கள்
உள்ள \Log{G2c} என்ற
முறைமையும் உண்டு
(காண்க~\olref{tab:G2c}).
\Log{G2c} இல் சுருக்கல்,
\Log{G1c} இல் இருப்பதைவிடவும்
முக்கியமானது.

\subfile{rules-G2c}

\end{document}
